\chapterhead{62}{ORR-23}{High-Level Summary of Basel III Reforms}{1}{2}

\lohead{62}{a}{Explain the motivations for revising the Basel III framework and
the goals and impacts of the December 2017 reforms to the Basel III framework.}

The goals and impacts of the December 2017 Basel III reforms include:

\begin{lettered}
\item Expanding the robustness and sensitivity of the standardised approach (SA)
for measuring credit risk, credit valuation adjustment (CVA) risk, and
operational risk.
\item Restricting the use of internal model approaches for credit risk, CVA risk,
and operational risk.
\item Introducing a leverage ratio buffer for global systemically important banks
(G-SIBs).
\item Creating an output floor that is more robust and risk-sensitive than the
current Basel II floor.
\end{lettered}

% ---------------------------------------------------------------------
\lohead{62}{b}{Summarize the December 2017 revisions to the Basel III framework in
the following areas: the standardized approach to credit risk; the internal
ratings-based (IRB) approaches for credit risk; the CVA risk framework; the
operational risk framework; and the leverage ratio framework.}

\orsubb{1) Standardized Approach (SA) for credit risk}
\begin{lettered}
\item Improved risk sensitivity by increasing granularity in risk-weight
definitions (e.g., residential mortgages vary by loan-to-value ratio).
\item Reduced reliance on external credit ratings.
\item Established the foundation for the revised output floor.
\end{lettered}

\orsubb{2) Internal Ratings-Based (IRB) approaches for credit risk}
\begin{lettered}
\item Restricted the advanced IRB (A-IRB) approach for large corporates and
financial institutions, mandating the foundation IRB (F-IRB) approach instead.
\item Introduced input floors for probability of default (PD), loss given default
(LGD), and exposure at default (EAD) to promote conservative risk estimates.
\end{lettered}

\orsubb{3) Credit Valuation Adjustment (CVA) risk framework}
\begin{lettered}
\item Enhanced risk sensitivity by including market-risk exposure and hedge
instruments.
\item Removed the IRB approach for CVA, requiring standardised or basic
approaches.
\item Aligned the CVA framework with the market risk framework, based on fair
value sensitivities.
\end{lettered}

\orsubb{4) Operational risk framework}
\begin{lettered}
\item Simplified the approach to a single Standardized Approach (SA) for all
banks.
\item Determined operational risk capital based on the bank's income and
historical operational risk losses.
\end{lettered}

\orsubb{5) Leverage ratio framework}
\begin{lettered}
\item Introduced a leverage ratio buffer specifically for Global Systemically
Important Banks (G-SIBs) that must be met with Tier 1 capital.
\item Refined measurements to better capture exposure from derivatives and
off-balance sheet items.
\end{lettered}

% ---------------------------------------------------------------------
\lohead{62}{c}{Describe the revised output floor introduced as part of the Basel
III reforms and approaches to be used when calculating the output floor.}

Before the crisis, banks using internal models could drastically lower capital
requirements compared to the standardised approach, gaining an unfair advantage.
Hence, the revised output floor sets a minimum capital requirement for internal
models.

\begin{center}
\begin{tikzpicture}[font=\fontsize{8.4}{10.5}\selectfont,
  bx/.style={text width=6.4cm,align=center,rounded corners=3pt,inner sep=9pt,
             text=navy,draw=palenavy,line width=0.7pt}]
\node[fill=navy,text=white,rounded corners=3pt,inner sep=6pt,text width=13.6cm,
      align=center,font=\sffamily\bfseries\fontsize{9}{11}\selectfont] (top) at (0,0)
      {A bank's risk-weighted assets (RWAs) must be the \emph{higher} of};
\node[bx,fill=paleblue,anchor=north east] (l) at ($(top.south east)+(-0.4,-0.5)$)
      {\textbf{72.5\%} of total risk-weighted assets using the SA};
\node[bx,fill=palerust,anchor=north west] (r) at ($(top.south west)+(0.4,-0.5)$)
      {RWAs from their \textbf{approved approach}\\[2pt]
       {\fontsize{7.6}{9.4}\selectfont(including both internal models and the
        standardised approach)}};
\node[font=\sffamily\bfseries\fontsize{9}{11}\selectfont,text=rust]
      at ($(l.west)!0.5!(r.east)$) {or};
\end{tikzpicture}
\end{center}
\orfigcap{The revised output floor, stated as a maximum of two measures.}

\orsub{Calculating the floor}

\noindent\begin{minipage}{\textwidth}
\renewcommand{\arraystretch}{1.35}
\begin{tabularx}{\textwidth}{@{}L{0.28\textwidth}Y@{}}
\rowcolor{navy} \thd{Risk type} & \thd{Approach used}\\
{\tabfont \textbf{1) Credit risk}} &
{\tabfont Standardised approach for credit risk (SA) is used.}\\
\rowcolor{paleblue}
{\tabfont \textbf{2) Counterparty credit risk}} &
{\tabfont Standardised approach for counterparty credit risk (SA-CCR) is used for
derivatives.}\\
{\tabfont \textbf{3) CVA risk}} &
{\tabfont Three options: SA-CVA, Basic Approach-CVA, or 100\% of the counterparty
credit risk capital requirement.}\\
\rowcolor{paleblue}
{\tabfont \textbf{4) Securitisation risk}} &
{\tabfont Three options: External Ratings Based Approach (SEC-ERBA), SEC-SA, or a
1,250\% risk weight.}\\
{\tabfont \textbf{5) Market and operational risk}} &
{\tabfont Standardised approach (SA) allowed for both.}\\
\end{tabularx}
\end{minipage}
\orfigcap{Which approach feeds the output floor, by risk type.}
